\documentclass[12pt,letterpaper]{article}

% Packages
\usepackage[margin=1in]{geometry}
\usepackage{amsmath,amssymb}
\usepackage{booktabs}
\usepackage{enumitem}
\usepackage{graphicx}
\usepackage{hyperref}
\usepackage{titlesec}
\usepackage{parskip}
\usepackage[T1]{fontenc}
\usepackage{lmodern}

% Formatting
\hypersetup{colorlinks=true, linkcolor=blue, urlcolor=blue, citecolor=blue}
\titleformat{\section}{\large\bfseries}{Section \thesection:}{0.5em}{}
\titleformat{\subsection}{\normalsize\bfseries}{}{0em}{}

% Title
\title{%
  \textbf{ISE 572: Industry 4.0 Machine Learning} \\[0.3em]
  \Large Final Project Proposal \\[0.5em]
  \large Machine Learning for Shipyard Production Scheduling: \\
  Neural Network-Based Block Sequencing and Processing Time Prediction
}
\author{Stephen Eacuello \and Steven Blum}
\date{March 13, 2026}

\begin{document}
\maketitle
\thispagestyle{empty}

% ============================================================
\section{Problem Statement}
% ============================================================

\subsection{The Industry 4.0 Problem}

Modern shipyards such as Hyundai Heavy Industries (HHI) Ulsan coordinate hundreds of steel blocks through a multi-stage production pipeline---steel cutting, panel assembly, block assembly, outfitting, painting, and erection at the dry dock. Blocks are moved by self-propelled modular transporters (SPMTs) and lifted into position by goliath cranes, all while respecting structural precedence constraints, equipment availability, and ship delivery deadlines. This makes shipyard scheduling a natural \textbf{process optimization} problem for Industry~4.0 methods.

\subsection{Why This Problem Matters}

Poor scheduling decisions cascade through the entire production line. An idle heavy crane can cost upwards of \$35,000 per day in lost throughput at a container terminal~\cite{gantrex2017}; shipyard goliath cranes with up to 900-ton lift capacity are in a comparable cost bracket. Late block delivery to the dock delays ship launches, triggering contractual penalties that can reach millions of dollars per vessel. Equipment failures from deferred maintenance cause unplanned downtime that compounds these costs further. Today, most yards rely on experienced human planners using heuristic rules and spreadsheets, which struggles to scale when a yard is handling 5--10 concurrent ship builds with 200--400 blocks each. Inefficient scheduling also increases energy consumption and creates downstream bottlenecks in outfitting and painting that grow over weeks.

\subsection{Industry Domain}

This project falls under \textbf{heavy manufacturing and maritime logistics}, specifically shipbuilding production management. The underlying problem structure---flexible job-shop scheduling with heterogeneous equipment and stochastic processing times~\cite{pernas2023}---also appears in aerospace assembly and modular construction.

% ============================================================
\section{Data Considerations}
% ============================================================

\subsection{Required Data}

The model needs five categories of data, all available within our simulation framework:

\textbf{Equipment telemetry and health data.} Real-time position, operational status (idle, transporting, lifting, under maintenance), and per-component health indicators for each SPMT (hydraulic system, tires, engine) and goliath crane (cable, motor). In a real yard this would come from IoT sensors and SCADA systems; our simulation models 24~SPMTs and 9~goliath cranes with a Wiener-process degradation model whose drift depends on load. We also maintain a prognostics and health management (PHM) module that estimates remaining useful life (RUL) using Kalman filtering, particle filtering, and Bayesian neural networks with uncertainty quantification---providing the scheduler with confidence intervals on when each piece of equipment is likely to fail.

\textbf{Block production records.} For each of the $\sim$50--200 blocks per ship: current production stage, processing times, assigned facility, weight, dimensions, and structural precedence relationships. In practice this comes from the yard's MES and ERP databases. Our environment tracks these as structured entities with full stage progression.

\textbf{Plate-level decomposition geometry.} Each block is composed of 10--50 steel plates characterized by dimensions (length, width, thickness), plate type (flat, curved, stiffened, bracket, bulkhead, shell), stiffener count, curvature radius, and material grade. This data originates from 3D CAD/CAM software and is exported through our decomposition scripts as structured JSON. These geometric features drive per-stage processing time predictions---a curved, stiffened plate takes much longer to cut and weld than a flat one.

\textbf{Industry benchmarking data.} We calibrate processing times against published NSRP and OECD data rather than inventing them. The OECD CGT system~\cite{oecd2007} provides ship-type complexity coefficients ($\text{CGT} = A \times \text{GT}^B$; for LNG carriers $A{=}32$, $B{=}0.68$), and the FMI Global Shipbuilding Benchmarking Study~\cite{fmi2005} reports man-hours per CGT across Korean, Japanese, European, and U.S.\ yards. Combined with stage-level breakdowns from Storch et al.~\cite{storch1995} and crew sizes from industry practice, these sources let us derive equipment-hours per block per stage for a 174,000\,cbm LNG carrier. We have compiled these into a structured benchmarks file and used them to fit regression coefficients across 3,000+ simulated block completions covering 7~production stages, achieving $R^2$ of 0.80--0.95 per stage.

\textbf{Ship schedule and constraint data.} Delivery deadlines, dock assignments, and the mapping between blocks and ships. In our simulation, up to 8~concurrent ships have staggered start dates and contract-derived due dates; when a ship is delivered, a new build is spawned to maintain yard load. Dock assignments respect physical dimensions (10~dry docks ranging from 260\,m to 672\,m in length) and crane reach. In practice this data comes from project management and contract databases.

\subsection{Data Challenges}

\textbf{Sim-to-real gap.} We train on simulated data with synthetic plate geometries, whereas real yards have noisy, heterogeneous data sources. To bridge this gap we use a calibration pipeline that fits regression coefficients from observed processing times, plus domain randomization that varies processing parameters during training.

\textbf{Plate data availability.} Full plate-level decomposition requires CAD post-processing that many yards do not routinely do for scheduling. As a fallback we generate synthetic plate data by estimating plate counts and geometry from block weight and type, validated against published shipbuilding references.

\textbf{Man-hours to equipment-hours conversion.} Industry benchmarks~\cite{fmi2005,storch1995} report labor in man-hours, but our simulation schedules equipment-hours (i.e., how long a facility is occupied). Converting between the two requires crew-size assumptions per stage (8--20 workers depending on the operation), which introduces modeling error if the assumed crew does not match a given yard's practice.

\textbf{Sparse reward signal.} Block completions happen infrequently (roughly every 50--100 scheduling steps), making credit assignment hard. We address this with reward shaping and dense intermediate signals like transport completion and stage advancement.

% ============================================================
\section{Proposed ML Solution}
% ============================================================

\subsection{Problem Formulation}

We frame the project around \textbf{four interconnected ML problems} that map to techniques covered in the course:

\subsubsection*{1.\ \ Regression: Processing Time Prediction from Plate Geometry}

Given a block's plate-level characteristics we predict processing time (hours) at each production stage. This is a \textbf{multivariate linear regression} problem:
\begin{equation}
  t_{\text{stage}} = \beta_0 + \beta_1 \cdot n_{\text{plates}} + \beta_2 \cdot n_{\text{curved}} + \beta_3 \cdot n_{\text{stiffened}} + \beta_4 \cdot A_{\text{total}} + \beta_5 \cdot L_{\text{weld}}
  \label{eq:regression}
\end{equation}
where the features are plate count, curved plate count, stiffened plate count, total plate area ($\text{m}^2$), and total weld length (m). We fit coefficients per stage (steel cutting, panel assembly, block assembly, etc.) using least-squares with non-negativity constraints. The plate decomposition data from CAD models provides ground-truth features for fitting, replacing the usual approach of historical averages or lognormal distributions. We validate with $R^2$, RMSE, and MAE using cross-validation across blocks.

\subsubsection*{2.\ \ Artificial Neural Networks: Scheduling Policy and Value Estimation}

The core scheduling agent uses a deep \textbf{actor-critic neural network} architecture:

\begin{itemize}[leftmargin=1.5em]
  \item \textbf{Actor network (policy):} A multi-layer perceptron (MLP) with ReLU activations that takes a learned state representation and outputs probability distributions over discrete actions. The network has a shared 2-layer trunk (256 hidden units) that branches into multiple classification heads: action type selection (4-way: dispatch SPMT, dispatch crane, schedule maintenance, or hold), equipment selection (which SPMT or crane), and target selection (which block/request to service). Each head applies softmax activation to produce a categorical probability distribution---this is a direct generalization of \textbf{logistic regression} to the multi-class setting, where a single-output sigmoid (binary logistic regression) extends to a $K$-way softmax over action categories. Additionally, the agent includes a binary \textbf{logistic regression classifier} (single sigmoid output) that predicts whether a block is ready to advance to its next production stage, used to gate dispatching decisions. Action masking sets infeasible actions to zero probability before sampling.

  \item \textbf{Critic network (value function):} A separate MLP head that estimates the expected future reward from the current state, used to compute advantage estimates for stable policy gradient training. The critic outputs a single scalar value via a 2-layer network ($256 \to 128 \to 1$).
\end{itemize}

The policy is trained using Proximal Policy Optimization (PPO) with clipped surrogate objectives, and alternatively with Soft Actor-Critic (SAC) using separate Q-networks for value estimation. We also employ DAgger (Dataset Aggregation), an imitation learning technique where the neural network learns from an expert scheduling heuristic and iteratively improves by collecting on-policy data.

\subsubsection*{3.\ \ Graph Neural Networks: Relational State Encoding (CNN Analog)}

The shipyard state has natural graph structure: blocks connect to facilities, SPMTs connect to blocks they can transport, cranes connect to blocks they can lift, and blocks have precedence relationships with each other. While a shipyard \emph{could} be laid out as a 2D spatial grid for \textbf{CNN} feature extraction~\cite{wang2021cnn}, key relationships like structural precedence and equipment eligibility are non-spatial and poorly served by fixed convolutional kernels. We instead use a \textbf{heterogeneous graph neural network (GNN)} with Graph Attention (GAT) layers, which generalizes the CNN's local aggregation to irregular, relational topologies~\cite{bronstein2021geometric}. Our architecture uses:

\begin{itemize}[leftmargin=1.5em]
  \item \textbf{Input projection layers:} Separate linear encodings per node type (blocks: 16~features including 4~plate-derived features, SPMTs: 10~features, cranes: 7~features, facilities: 3~features).
  \item \textbf{Message-passing layers:} Two layers of multi-head graph attention convolution (4~attention heads per layer) with residual connections and layer normalization, operating over 8~edge types (e.g., ``block needs\_transport SPMT'', ``block precedes block'').
  \item \textbf{Global pooling:} Mean pooling per node type, concatenated into a fixed-size state embedding (512~dimensions) that feeds the actor-critic ANNs described above.
\end{itemize}

Four plate-derived features---normalized plate count, normalized plate area, percent curved plates, and percent stiffened plates---give the GNN richer block representations that capture manufacturing complexity beyond simple weight and size.

\subsubsection*{4.\ \ Bayesian Inference: Equipment Remaining Useful Life Estimation}

The scheduling agent needs to decide \emph{when} to pull equipment offline for maintenance. We frame this as a \textbf{Bayesian state estimation} problem: given a sequence of noisy health observations from the Wiener degradation model, estimate the posterior distribution over remaining useful life (RUL) before failure. Our PHM module implements three estimators of increasing flexibility:

\begin{itemize}[leftmargin=1.5em]
  \item \textbf{Kalman filter:} Closed-form Gaussian posterior for linear degradation dynamics. Fast but assumes constant drift.
  \item \textbf{Particle filter:} Sequential Monte Carlo for nonlinear, load-dependent degradation. Captures asymmetric failure distributions.
  \item \textbf{Bayesian neural network:} Data-driven estimator that outputs mean, standard deviation, and 95\% confidence intervals on RUL, trained on simulated degradation trajectories.
\end{itemize}

The RUL estimates feed directly into the scheduling policy: the agent observes RUL confidence intervals as part of the equipment state and learns to schedule preventive maintenance before the lower bound approaches zero, avoiding unplanned breakdowns.

\subsection{Stochastic Modeling Extensions}

Beyond the core ML problems above, we are investigating how several real-world stochastic factors affect scheduling performance. These are not separate ML models but rather extensions to the simulation environment that test policy robustness:

\begin{itemize}[leftmargin=1.5em]
  \item \textbf{Spatial constraints:} Dock and staging area capacity limits, crane reach envelopes, and SPMT routing conflicts that restrict which blocks can be worked on simultaneously.
  \item \textbf{Labor resource leveling:} Crew availability varies across shifts and trades; overbooking a trade (e.g., welders) at one facility starves another.
  \item \textbf{Duration uncertainty:} Processing times are drawn from calibrated distributions rather than fixed means, so the policy must be robust to blocks finishing earlier or later than predicted.
  \item \textbf{Shift start/stop patterns:} Discrete work shifts introduce idle gaps and batch-start effects that continuous-time models miss.
  \item \textbf{Weather disruptions:} Outdoor operations (crane lifts, painting) can be delayed by wind or rain events sampled from seasonal distributions.
\end{itemize}

We evaluate the trained policy's sensitivity to each factor by toggling it on/off and measuring the change in throughput and tardiness, which tells us which stochastic elements matter most for scheduling quality.

\subsection{Baseline Comparisons}

We compare against traditional operations research baselines:

\begin{itemize}[leftmargin=1.5em]
  \item \textbf{Earliest Due Date (EDD):} Rule-based heuristic (no ML)
  \item \textbf{Mixed-Integer Programming (MIP):} PuLP CBC solver with rolling-horizon optimization
  \item \textbf{Genetic Algorithm:} Evolutionary optimization with permutation chromosomes
  \item \textbf{Model Predictive Control:} CP-SAT constraint programming
\end{itemize}

% ============================================================
\section{Expected Impact}
% ============================================================

\subsection{Industry Benefits}

A scheduling policy that outperforms expert heuristics would reduce block throughput time, equipment idle time, and delivery tardiness. In our preliminary experiments the Expert EDD baseline erects $\sim$32 blocks per episode (300~time steps) while MIP achieves $\sim$25. An RL policy that matches or beats Expert performance while also timing maintenance better would provide:

\begin{itemize}[leftmargin=1.5em]
  \item \textbf{Reduced idle costs:} Fewer idle crane-hours through better sequencing.
  \item \textbf{Faster throughput:} More blocks completed per unit time by optimizing SPMT routing and crane assignment jointly.
  \item \textbf{Predictive maintenance integration:} Scheduling maintenance during natural lulls rather than fixed intervals, reducing unplanned downtime.
  \item \textbf{Accurate time estimates:} Plate-geometry-driven processing time regression replacing rough historical averages, enabling tighter schedule planning.
\end{itemize}

Even without the RL policy, the plate-level processing time model is useful on its own---per-stage estimates based on actual geometry are more informative than yard-wide averages for capacity planning.

\subsection{Evaluation Metrics}

We evaluate using standard scheduling and shipyard operations metrics:

\begin{table}[h]
\centering
\small
\begin{tabular}{@{}p{3.2cm}p{6.5cm}p{4.2cm}@{}}
\toprule
\textbf{Metric} & \textbf{Definition} & \textbf{Target} \\
\midrule
Throughput & Blocks erected per time unit & Maximize ($\geq$ Expert) \\
Total tardiness & $\sum(\text{completion} - \text{due date})$ for late blocks & Minimize \\
Equip.\ utilization & Fraction of time equipment is productively busy & $> 70\%$ \\
Makespan & Time to complete all blocks for a ship & Minimize \\
Proc.\ time $R^2$ & Regression accuracy of plate-based predictions & $> 0.80$ (prelim.\ 0.80--0.95) \\
RMSE & Root mean squared error of time predictions & $< 2$ hours/stage \\
\bottomrule
\end{tabular}
\caption{Evaluation metrics and targets.}
\label{tab:metrics}
\end{table}

We also report wall-clock time per decision to check that the policy is fast enough for real-time use (target $< 10$\,ms per action, vs.\ seconds for MIP solvers).

% ============================================================
% References
% ============================================================

\begin{thebibliography}{99}

% --- Cited in text, in order of appearance ---

\bibitem{gantrex2017}
{Gantrex},
``The cost of crane rail quality: How not to lose {USD}~35,000 a day,''
Blog post, March 2017. Available: \url{https://www.gantrex.com/the-cost-of-crane-rail-quality-how-not-to-lose-usd-35000-a-day/}.

\bibitem{pernas2023}
J.~Pernas-\'{A}lvarez and D.~Crespo-Pereira,
``A constrained programming model for the optimization of industrial-scale scheduling problems in the shipbuilding industry,''
\textit{Journal of Marine Science and Engineering}, vol.~11, no.~8, p.~1517, 2023.

\bibitem{oecd2007}
{OECD Council Working Party on Shipbuilding},
``Compensated gross ton ({CGT}) system,''
OECD Document C/WP6(2006)7, 2007.

\bibitem{fmi2005}
{First Marine International},
``Global shipbuilding industrial base benchmarking study---Part~1: Major shipyards,''
Report for the National Shipbuilding Research Program ({NSRP}), May 2005. Available: \url{https://www.nsrp.org/wp-content/uploads/2024/07/FMI-Global_Industrial_Benchmarking-major_yards-Navy-rpt.pdf}.

\bibitem{storch1995}
R.~L.~Storch, C.~P.~Hammon, H.~M.~Bunch, and R.~C.~Moore,
\textit{Ship Production}, 2nd~ed.
Cornell Maritime Press, 1995.

\bibitem{wang2021cnn}
L.~Wang, Z.~Pan, and J.~Wang,
``A review of reinforcement learning based intelligent optimization for manufacturing scheduling,''
\textit{Complex System Modeling and Simulation}, vol.~1, no.~4, pp.~257--270, 2021.

\bibitem{bronstein2021geometric}
M.~M.~Bronstein, J.~Bruna, T.~Cohen, and P.~Veli\v{c}kovi\'{c},
``Geometric deep learning: Grids, groups, graphs, geodesics, and gauges,''
\textit{arXiv preprint arXiv:2104.13478}, 2021.

% --- Additional references ---

\bibitem{schulman2017ppo}
J.~Schulman, F.~Wolski, P.~Dhariwal, A.~Radford, and O.~Klimov,
``Proximal policy optimization algorithms,''
\textit{arXiv preprint arXiv:1707.06347}, 2017.

\bibitem{ross2011dagger}
S.~Ross, G.~J.~Gordon, and D.~Bagnell,
``A reduction of imitation learning and structured prediction to no-regret online learning,''
in \textit{Proc.\ AISTATS}, pp.~627--635, 2011.

\bibitem{zhou2023}
T.~Zhou, L.~Luo, Y.~He, Z.~Fan, and S.~Ji,
``Solving panel block assembly line scheduling problem via a novel deep reinforcement learning approach,''
\textit{Applied Sciences}, vol.~13, no.~14, p.~8483, 2023.

\bibitem{lamb2004}
T.~Lamb, Ed.,
\textit{Ship Design and Construction}, vols.~1--2.
Society of Naval Architects and Marine Engineers ({SNAME}), 2003--2004.

\end{thebibliography}

\end{document}
